% Notes of the Montreal seminar, 1970 (cote 129, typed part): archive cote 129, site page 20 (PDF page 21).
% typescript. Draft compared: none.
% AI-transcribed by Claude Opus 5.5 (claude-opus-5-5, adaptive thinking, effort high), Message Batches API, batch mtl129_full (msgbatch_01AGnDKHwy6ATCM7ehUa6263), 2026-11-13; not yet checked by eye.
\scanpage{129/20}
phe à $\mathbb{E}^{(\mathbb{N})}$ , c'est-à-dire que pour tout anneau $A$ ,
\[
\underset{\to}{\mathbb{W}}(A) = A^{(\mathbb{N})} = \{(x_1,x_2,\ldots,x_n,\ldots) \ , \ x_i \in A \ , \ x_i = 0 \ \text{ sauf un nombre fini}\} \ .
\]
Mais ce foncteur a une structure additive qui n'est pas la structure
habituelle de $\mathbb{E}^{(\mathbb{N})}$ .

On peut définir un autre système inductif naturel à l'aide
des morphismes
\[
\mathbb{W}_n \xrightarrow{T_n F_n} \mathbb{W}_{n+1} \ .
\]

On a alors des diagrammes commutatifs
\diagram{
\begin{tikzcd}
\mathbb{W}_n \arrow[r, "T_n"] \arrow[d, "F_n^{n-1}"] & \mathbb{W}_{n+1} \arrow[d, "F_{n+1}^{n}"] \\
\mathbb{W}_n \arrow[r, "T_n F_n = p"] & \mathbb{W}_{n+1}
\end{tikzcd}
} ,

car $T_n F_n = F_{n+1} T_n$ . Par définition, on pose
\[
\mathbb{W}' = \underset{\to}{\lim}(\mathbb{W}_n, T_n F_n) \ .
\]

En caractéristique $0$ , $\underset{\to}{\mathbb{W}}'$ n'a aucune structure naturelle de Groupe,
mais en caractéristique $p$ $\underset{\to}{\mathbb{W}}'$ est un ind-schéma en groupes, et on a un
morphisme additif
\[
u : \underset{\to}{\mathbb{W}}_{\mathbb{F}_p} \to \underset{\to}{\mathbb{W}}'_{\mathbb{F}_p}
\]
obtenu par passage à la limite des $F^{n-1}$ .

\mspage{\hand{20}}
